% Generated draft: requires mathematical and source review.
\documentclass[11pt]{article}
\usepackage[margin=25mm]{geometry}
\usepackage{fontspec}
\setmainfont{Latin Modern Roman}
\usepackage{amsmath,amssymb,mathtools,unicode-math}
\setmathfont{Latin Modern Math}
\usepackage{enumitem,needspace,xurl,hyperref}
\setlength{\parindent}{0pt}
\setlength{\parskip}{5pt}
\setlength{\emergencystretch}{3em}
\newcommand{\N}{\mathbb{N}}
\newcommand{\Z}{\mathbb{Z}}
\newcommand{\Q}{\mathbb{Q}}
\newcommand{\R}{\mathbb{R}}
\newcommand{\C}{\mathbb{C}}
\newcommand{\F}{\mathbb{F}}
\newcommand{\A}{\mathbb{A}}
\newcommand{\PP}{\mathbb{P}}
\newcommand{\E}{\mathbb{E}}
\newcommand{\eps}{\varepsilon}
\newcommand{\dd}{\,\mathrm{d}}
\begin{document}
\section*{UnsolvedMath ID: 30003558}
\textbf{Source problem code:} catalogue entry
\subsection{Definitions and mathematical statement}
The supplied record names the Strong Slope Conjecture but does not specify its degree normalization or numerical identities. The following is a conventional formulation for knots in \(S^3\); these conventions and their agreement with the cited report require human verification.

Let \(K\subset S^3\) be a knot and let \(M_K=S^3\setminus\operatorname{int}\nu(K)\) be its exterior. Fix the meridian \(\mu\) and preferred longitude \(\lambda\). A boundary slope \(r/s\in\Q\), written with \(\gcd(r,s)=1\) and \(s>0\), means that each boundary component represents the unoriented primitive class \(r\mu+s\lambda\). An essential surface here means a compact, properly embedded surface with nonempty boundary that is incompressible and boundary-incompressible and has no boundary-parallel component. For such a surface with boundary slope \(r/s\), its number of sheets is \(s|\partial S|\), where \(|\partial S|\) is its number of boundary components.

Use the unreduced colored Jones polynomial \(J_K(n;q)\), with positive integer color \(n\) and unknot normalization
\[
J_U(n;q)=\frac{q^{n/2}-q^{-n/2}}{q^{1/2}-q^{-1/2}}.
\]
Write its maximal and minimal \(q\)-degrees, for sufficiently large \(n\), as quadratic quasi-polynomials
\[
d_+(n)=a_+(n)n^2+b_+(n)n+c_+(n),\qquad
d_-(n)=a_-(n)n^2+b_-(n)n+c_-(n),
\]
where the coefficient functions are periodic. For a real sequence \((x_n)\), let \(\operatorname{Acc}(x_n)\) denote its set of subsequential limits as \(n\to\infty\). Define
\[
\begin{aligned}
\mathcal J_+(K)&=\operatorname{Acc}\!\left(4d_+(n)/n^2\right),&
\mathcal X_+(K)&=\operatorname{Acc}\!\left(2b_+(n)\right),\\
\mathcal J_-(K)&=\operatorname{Acc}\!\left(4d_-(n)/n^2\right),&
\mathcal X_-(K)&=\operatorname{Acc}\!\left(2b_-(n)\right).
\end{aligned}
\]
The conventional Strong Slope Conjecture asks whether, for every knot \(K\subset S^3\), both of the following hold:
\begin{enumerate}
\item For every \(r/s\in\mathcal J_+(K)\) in reduced form with \(s>0\), there exists an essential surface \(S_+\subset M_K\) with boundary slope \(r/s\) such that
\[
\frac{\chi(S_+)}{s|\partial S_+|}\in\mathcal X_+(K).
\]
\item For every \(r/s\in\mathcal J_-(K)\) in reduced form with \(s>0\), there exists an essential surface \(S_-\subset M_K\) with boundary slope \(r/s\) such that
\[
-\frac{\chi(S_-)}{s|\partial S_-|}\in\mathcal X_-(K).
\]
\end{enumerate}
Surfaces satisfying the corresponding conditions are called Jones surfaces. This formulation asserts existence of surfaces realizing degree-derived slopes and normalized Euler characteristics. It does not assert that every boundary slope is detected, or that Euler characteristic and number of sheets are separately or uniquely determined. The supplied wording about determining these quantities therefore needs qualification. The record also does not establish whether a stronger correspondence between coefficients along the same residue class is intended.
\subsection{Short English statement}
For every knot, the Strong Slope Conjecture predicts essential surfaces whose boundary slopes come from quadratic colored Jones degree growth and whose Euler characteristics divided by their numbers of sheets match linear degree data. The displayed formulation uses explicit conventional choices that require comparison with the original source; it does not claim separate determination of Euler characteristics or sheet counts.
\subsection{Sources}
\begin{itemize}
\item Supplied local registry record, permanent UnsolvedMath ID \(30003558\), source problem code \texttt{catalogue entry}, title: \emph{The Strong Slope Conjecture for Colored Jones Polynomials}. The record belongs to the Oberwolfach Reports Open Problems collection and gives proposed year 2017; no proposer is identified.
\item Supplied source citation: \emph{Low-dimensional Topology and Number Theory} (2018), Oberwolfach Reports, pp.\ 2363--2426, DOI \texttt{10.4171/owr/2017/38}. \url{https://doi.org/10.4171/owr/2017/38}. The background describes Oberwolfach Reports as documenting workshops and their open-problem discussions. The supplied publication year and the year embedded in the DOI are retained without reconciliation.
\item Supplied literature source: \emph{Slope conjectures for 3 string Montesinos knots}. \url{https://nyjm.albany.edu/j/2019/25-2.html}.
\item Supplied literature source: \emph{The Slope Conjecture for Montesinos knots}. \url{https://arxiv.org/abs/1807.00957}.
\item Literature assessment dated \texttt{2026-08-21}: the supplied curation labels the problem partially solved, reports proofs for important classes including broad families of Montesinos knots, and reports no general proof. This is supplied dated curation derived from the OWR record, not an independently verified assessment; its time-sensitive status and cited claims require independent verification.
\item The registry labels the statement \texttt{corrected\_verified} and says its mathematical content was preserved while notation, typography, or omitted definitions were repaired. This is supplied registry metadata, not mathematical verification of this draft. The cited sources have not been independently consulted here.
\end{itemize}
\end{document}
